% !TeX root = ../../main.tex
\section{The top Hodge class and Bessel asymptotics}
\label{sec:s06}

The top-class restriction in
Corollary~\ref{crl:s02:defect-allocation} removes every graph with a
cycle and factors $\lambda_g$ over the remaining trees. Combined
with the boundary exponential~\eqref{eq:s02:boundary-exponential},
it turns the integral defining $b_g$
into a sum of top-Hodge vertex integrals. This gives a proof of
Theorem~\ref{thm:top-hodge} using the geometry of
Section~\ref{sec:s02} and classical Hodge and Bessel formulas.

There are four steps. The $\lambda_g$ theorem first reduces each
vertex to a one-variable coefficient. A strict modulus estimate on
a Bessel circle then bounds all stable vertex types at once. This
uniform estimate controls the entire tree tail, leaving only the
tree with no edges and an elliptic leaf at the precision of the theorem.
Finally a two-term saddle calculation and the leaf integral give
the three contributions to the scalar $b$. The exact genus-zero
reduction, Bessel inverse and critical-point data are retained for
the later fixed-defect analysis.

% !TeX root = ../../main.tex
\subsection{Top Hodge intersections and exact vertices}
\label{subsec:s06:01}

We use Faber--Pandharipande's multipoint theorem
\cite[Theorem~1]{FP03}, together with their earlier one-point
calculation \cite[Theorem~2, (6)]{FP00}. Set $\beta_0=1$ by
convention and, for an integer $G\geq1$, define
\begin{equation}\label{eq:s06:beta-definition}
 \beta_G=\int_{\MM_{G,1}}\lambda_G\psi_1^{2G-2}.
\end{equation}
In particular $\beta_0$ is not an integral on a genus-zero one-point
space.

\begin{thm}[Top Hodge intersections]\label{thm:s06:lambda-g}
For integers $G\geq1$, $N\geq1$ and nonnegative integers
$d_1,\ldots,d_N$ with $\sum_i d_i=2G-3+N$,
\begin{equation}\label{eq:s06:lambda-g}
 \int_{\MM_{G,N}}\lambda_G\prod_{i=1}^N\psi_i^{d_i}
 =\beta_G\frac{(2G-3+N)!}{\prod_{i=1}^N d_i!}.
\end{equation}
The same formula holds for $G=0$, $N\geq3$ with $\beta_0=1$.
For a formal variable $z$,
\begin{equation}\label{eq:s06:beta-series}
 \sum_{G\geq0}\beta_Gz^{2G}=\frac{z/2}{\sin(z/2)}.
\end{equation}
As $G\to\infty$,
\begin{equation}\label{eq:s06:beta-asymptotics}
 \beta_G=2(2\pi)^{-2G}(1+O(4^{-G})),\qquad
 \frac{\beta_{G-1}}{\beta_G}=(2\pi)^2(1+O(4^{-G})).
\end{equation}
There are absolute constants $c,C>0$ such that
$c(2\pi)^{-2G}\leq\beta_G\leq C(2\pi)^{-2G}$ for every $G\geq0$.
\end{thm}
\begin{prf}
The positive-genus formula is \cite[Theorem~1]{FP03}; the
stable genus-zero formula is its equation~(1). The generating
series is \cite[Theorem~2, (6)]{FP00}, whose genus-zero value is
one. For real $s>1$, write
$\zeta(s)=\sum_{k\geq1}k^{-s}$ for the Riemann zeta function.
The Bernoulli evaluation in \cite[\S 4.2, Lemma~1]{FP00}
and \cite[(25.6.2)]{DLMF} give, for $G\geq1$,
\[
 \beta_G=(1-2^{1-2G})\frac{|B_{2G}|}{(2G)!}
       =2(1-2^{1-2G})\zeta(2G)(2\pi)^{-2G}.
\]
The Dirichlet series \cite[(25.2.1)]{DLMF} implies
$\zeta(2G)=1+O(4^{-G})$ as $G\to\infty$.
This proves both asymptotics and the uniform comparison, extending
the constants to the finitely many small genera and $G=0$.
\end{prf}

For $x\in\mathbb C$, define the entire function
\begin{equation}\label{eq:s06:phi}
 \Phi(x)=\sum_{a\geq0}\frac{(-1)^a x^a}{a!(a+1)!}
        =\frac{J_1(2\sqrt x)}{\sqrt x}.
\end{equation}
The quotient and every subsequent square-root expression in Bessel
functions mean their entire power series.

For a stable pair $(G,n)$, let $\mathbf s=(s_1,\ldots,s_n)$
be a complex vector and set
\begin{equation}\label{eq:s06:vertex-definition}
 W^{\rm top}_{G,n}(\mathbf s)
  =\int_{\MM_{G,n}}\lambda_G
                   e^{\kappa_1-\sum_{i=1}^n s_i\psi_i}.
\end{equation}
The superscript distinguishes this vertex integral from other
intersection generating functions. The local coefficient order
$m=2G-3+n$ is nonnegative for every stable type.

The whole exponential specialization
\eqref{eq:s02:kappa-exponential} is especially useful here. It
replaces the background by new cotangent insertions, after which
\eqref{eq:s06:lambda-g} determines every coefficient.

\begin{prp}[Exact vertices]\label{prp:s06:exact-vertex}
For every stable pair $(G,n)$, with $m=2G-3+n$,
\begin{equation}\label{eq:s06:exact-vertex}
 W^{\rm top}_{G,n}(\mathbf s)
 =\beta_G m![x^m]\Phi(x)^{-m-1}
                       \prod_{i=1}^n J_0(2\sqrt{s_i x}).
\end{equation}
\end{prp}
\begin{prf}
Apply \eqref{eq:s02:kappa-exponential} with scalar parameter one.
A term with $k\geq0$ forgotten points has coefficient $1/k!$ and
new-point factor $\psi(1-e^{-\psi})$ at each point.
The Hodge insertion pulls back unchanged. The new powers have
minimum degree two, so Lemma~\ref{lmm:s02:rational-tail-degree}
replaces the old cotangent pullbacks by their intrinsic classes.
This is the inversion of the entire kappa background; there is no
further background exponential to pull back.

The multinomial formula contributes $\beta_G(m+k)!$ and a
factorial denominator for each descendant degree. Thus the old
marking $i$ contributes
\[
 \sum_{a\geq0}\frac{(-s_i x)^a}{(a!)^2}
       =J_0(2\sqrt{s_i x}),
\]
whereas a new marking contributes
\[
 \sum_{a\geq1}\frac{(-1)^{a-1}x^{a+1}}{a!(a+1)!}
       =x(1-\Phi(x)).
\]
After removing $x^k$ from the coefficient extraction, the vertex
integral is therefore
\[
 \beta_G[x^m]\prod_{i=1}^nJ_0(2\sqrt{s_i x})
       \sum_{k\geq0}\frac{(m+k)!}{k!}(1-\Phi(x))^k.
\]
This sum is coefficientwise finite because $1-\Phi(x)$ has zero
constant term. The binomial series makes it
$m!\Phi(x)^{-m-1}$, proving the formula.

For an unmarked vertex $G\geq2$, one has $m\geq1$. Its $k=0$
term vanishes by degree, and every nonzero term has at least one
marking, as required by the positive-genus multipoint theorem.
For $G=0$ the original stability already gives $n\geq3$.
These observations cover the ranges excluded from a literal
application of \eqref{eq:s06:lambda-g} with $N=0$.
\end{prf}

\begin{crl}[Numerical genus-zero reduction]\label{crl:s06:genus-zero}
For every stable type,
\begin{equation}\label{eq:s06:genus-zero}
 W^{\rm top}_{G,n}(\mathbf s)
 =\beta_G W^{\rm top}_{0,2G+n}(\mathbf s,0,\ldots,0).
\end{equation}
\end{crl}
\begin{prf}
The type on the right is stable since $2G+n=m+3\geq3$.
It has the same coefficient order $m$, and its additional $2G$
old-marking factors are $J_0(0)=1$ in
\eqref{eq:s06:exact-vertex}. Its one-point convention is $\beta_0=1$.
\end{prf}

\begin{rmk}[Classical inverse]
This is an equality of numerical generating functions on different
moduli spaces. The inverse of $x\mapsto x\Phi(x)$ is the classical
genus-zero kappa series of \cite[(0.8)]{KMZ96v1}, with its earlier
Zograf attribution. For a formal variable $t$ and an integer $N\geq3$,
Lagrange inversion and \eqref{eq:s06:exact-vertex} give its
$t^{N-2}$ coefficient as
\[
 \frac{W^{\rm top}_{0,N}(\mathbf0)}{(N-2)!}
 =\frac{\displaystyle\int_{\MM_{0,N}}\kappa_1^{N-3}}
        {(N-2)!(N-3)!}.
\]
This makes the conversion from raw kappa powers to exponential
integrals explicit. The inverse and critical behavior are also
analyzed in \cite[\S 6, proof of 6.3]{MaZ00}. The uniform vertex
and complete tree estimates needed here will be proved directly.
\end{rmk}

% !TeX root = ../../main.tex
\subsection{The Bessel critical point and uniform bounds}
\label{subsec:s06:02}

Using the zero fixed in the Introduction, define
\begin{equation}\label{eq:s06:critical-data}
 x_* =\frac{z_{0,1}^2}{4},\qquad
 w(x)=x\Phi(x),\qquad
 R=w(x_*)=\frac{z_{0,1}J_1(z_{0,1})}{2}=\frac{\rho}{12}.
\end{equation}
The strict circle estimate below makes this positive point the
unique dominant point of the Cauchy integral. It also supplies a
bound that is uniform when both the genus and the valence vary.

\begin{lmm}[Critical circle]\label{lmm:s06:critical-circle}
The function $\Phi$ has no zero on $|x|\leq x_*$, and
\begin{equation}\label{eq:s06:critical-derivatives}
 w'(x)=J_0(2\sqrt x),\quad w'(x_*)=0,\quad
 w''(x_*)=-\frac{2J_1(z_{0,1})}{z_{0,1}}<0.
\end{equation}
There is a constant $a>0$ such that, for $-\pi\leq\theta\leq\pi$,
\begin{equation}\label{eq:s06:circle-bound}
 |\Phi(x_*e^{i\theta})|
       \geq\Phi(x_*)e^{a\theta^2}.
\end{equation}
In particular $R>0$.
\end{lmm}
\begin{prf}
For an integer $k\geq1$, let $z_{1,k}$ be the $k$th positive zero
of $J_1$. The product \cite[(10.21.15)]{DLMF} gives
\begin{equation}\label{eq:s06:phi-product}
 \Phi(x)=\prod_{k\geq1}\left(1-\frac{4x}{z_{1,k}^2}\right).
\end{equation}
Interlacing \cite[(10.21.2)]{DLMF} gives
$z_{0,1}<z_{1,1}$. Thus all factors are positive at $x_*$ and all
zeros of $\Phi$ lie outside its closed disk.
The derivative identities \cite[(10.6.6), (10.6.3)]{DLMF} give
\eqref{eq:s06:critical-derivatives}. Also $J_1$ is positive between
zero and $z_{1,1}$, so $J_1(z_{0,1})>0$ and $R>0$.

For $0<u<1$ one has
\[
 |1-ue^{i\theta}|^2=(1-u)^2+2u(1-\cos\theta).
\]
Consequently each factor of \eqref{eq:s06:phi-product} has modulus
at least its value at $x_*$. For the first factor, the logarithm
of the ratio divided by $\theta^2$ extends to a strictly positive
continuous function at zero and is positive elsewhere on
$[-\pi,\pi]$. Its positive minimum gives $a$ in
\eqref{eq:s06:circle-bound}; the other factors only improve the
bound. The convergent product permits passage from finite products
to $\Phi$.
\end{prf}

\begin{rmk}[Critical-value convention]
The critical value denoted $x_0$ in
\cite[\S 6, (17) and proof of 6.3]{MaZ00} is our $R$.
Our positive critical argument $x_*$ and its derivative properties
are fixed by \eqref{eq:s06:critical-data} and proved above.
\end{rmk}

\begin{lmm}[All stable vertex types]\label{lmm:s06:vertex-bound}
There is a constant $C>0$, independent of the stable type $(G,n)$,
such that for $s_i\in[0,1]$ and $m=2G-3+n$,
\begin{equation}\label{eq:s06:vertex-bound}
 |W^{\rm top}_{G,n}(\mathbf s)|
 \leq C^{n+1}\beta_G\frac{m!R^{-m}}{\sqrt{m+1}}.
\end{equation}
\end{lmm}
\begin{prf}
On $|x|=x_*$, the entire series for $J_0$ gives the bound
\[
 |J_0(2\sqrt{s x})|
 \leq\sum_{k\geq0}\frac{x_*^k}{(k!)^2}
 \qquad(0\leq s\leq1).
\]
Its right side is a fixed finite constant. Cauchy's formula and
\eqref{eq:s06:circle-bound} bound the coefficient in
\eqref{eq:s06:exact-vertex} by
\[
 C^n x_*^{-m}\Phi(x_*)^{-m-1}
 \frac1{2\pi}\int_{-\pi}^{\pi}e^{-a(m+1)\theta^2}\dd\theta
 \leq\frac{C^{n+1}R^{-m}}{\sqrt{m+1}}.
\]
Multiplication by $\beta_Gm!$ proves the assertion, also for $m=0$.
The constants depend only on the fixed circle, not on $G$ or $n$.
\end{prf}

% !TeX root = ../../main.tex
\subsection{Stable trees and their tails}
\label{subsec:s06:03}

For a vertex $v$ of a tree, put $m_v=2g_v-3+n_v$. Let
$\mathbf s_v$ be its tuple of flag parameters: both flags belonging
to an edge $e$ use the same $s_e\in[0,1]$.
For the comparisons in this and the next subsection define the
positive scale
\begin{equation}\label{eq:s06:tree-scale}
 \mathcal T_g=\beta_g\frac{g_2!R^{-g_2}}{\sqrt{g_2+1}}.
\end{equation}
We first keep the exact automorphism factors. Only after bounding
every vertex will we replace unrooted trees by plane rooted trees
in a positive majorant. This controls the entire genus and valence
sum at each edge count.

\begin{prp}[The complete tree sum]\label{prp:s06:tree-tail}
For $g\geq2$,
\begin{equation}\label{eq:s06:tree-expansion}
 \frac{12^{g_2}b_g}{g_2!}
 =\sum_{\substack{\Gamma\in\mathcal G_{g,0}\\\Gamma\text{ a tree}}}
  \frac1{|\operatorname{Aut}\Gamma|}
  \int_{[0,1]^{E(\Gamma)}}
       \prod_{v\in V(\Gamma)}W^{\rm top}_{g_v,n_v}(\mathbf s_v)
       \prod_{e\in E(\Gamma)}\dd s_e.
\end{equation}
Every contributing tree has $E\leq g_2$ edges. There is a constant
$C>0$, independent of $g$ and $E$, such that the absolute sum of
its $E$-edge integrals is at most
\begin{equation}\label{eq:s06:edge-mass}
 \mathcal T_g C^{E+1}
 \sqrt{\frac{g_2+1}{g_2-E+1}}\,
                  \frac{(g_2-E)!}{g_2!}
 \qquad(0\leq E\leq g_2).
\end{equation}
Hence, for every fixed integer $K\geq1$, the sum over all trees
with at least $K$ edges is $O_K(\mathcal T_g g_2^{-K})$ in
absolute value.
\end{prp}
\begin{prf}
\emph{The finite-genus identity.}
The boundary exponential with both scalar coefficients one gives
edge kernel $\int_0^1e^{-s(\psi_f+\psi_{f'})}\dd s$.
By Corollary~\ref{crl:s02:defect-allocation}, the pullback of
$\lambda_g$ vanishes on every graph with a cycle and is the product
of the vertex top classes on a tree. Logarithmic kappa additivity
\eqref{eq:s02:tautological-gluing} now gives the right side of
\eqref{eq:s06:tree-expansion}, with the original graph
$1/|\operatorname{Aut}\Gamma|$.
On the left, $\kappa_1+\delta=12\lambda_1$ and degree conversion
gives $\int\lambda_ge^{12\lambda_1}=12^{g_2}b_g/g_2!$.
For a tree, the genus condition and flag count imply
\begin{equation}\label{eq:s06:tree-degree}
 \sum_v m_v=\sum_v(2g_v-3+n_v)=g_2-E.
\end{equation}
All $m_v$ are nonnegative, proving $E\leq g_2$.

\emph{A majorant for every tree and genus assignment.}
The comparison for $\beta_G$ in Theorem~\ref{thm:s06:lambda-g}
uses $\sum_v g_v=g$ to give
$\prod_v\beta_{g_v}\leq C^{E+1}\beta_g$, including rational and
elliptic vertices. Apply Lemma~\ref{lmm:s06:vertex-bound} to
every factor. Since $\sum_v n_v=2E$ and there are $E+1$ vertices,
all its exponential constants are bounded by $C^{E+1}$ after
increasing $C$. The edge domains have measure one. After removing
$\beta_gR^{-g_2}$, the remaining majorant is thus a constant to
power $E+1$ times
$\prod_v m_v!/\sqrt{m_v+1}$.

Every unrooted genus-labelled tree has a plane rooted realization,
and its original weight is at most one. A plane rooted tree with
$E$ edges is encoded by a walk of $2E$ up/down steps, so there are
at most $4^E$ shapes. For a fixed shape and the $m_v$, the genus
labels are determined by $2g_v=m_v+3-n_v$. The beta comparison has
already removed them; dropping parity and stability restrictions
therefore only enlarges the positive sum over nonnegative $m_v$.
This introduces no factorial in the vertex count.

For every integer $M\geq0$,
\begin{equation}\label{eq:s06:factorial-convolution}
 \sum_{k=0}^M k!(M-k)!\leq3M!.
\end{equation}
Indeed the two endpoint terms give $2M!$ when $M\geq1$, while
$\binom{M}{k}\geq M$ for $1\leq k\leq M-1$; the case $M=0$
is immediate. Iterating the convolution for $E+1$ vertices gives
$3^EM!$ for total $M=\sum_v m_v$. Also
$\prod_v(m_v+1)\geq M+1$. With $M=g_2-E$, these estimates and
the shape count prove \eqref{eq:s06:edge-mass}.

\emph{Summing every edge count.}
For $0\leq E\leq g_2$, direct multiplication gives
\[
 (E+1)(g_2-E+1)-(g_2+1)=E(g_2-E)\geq0.
\]
The square root in \eqref{eq:s06:edge-mass} is therefore at most
$\sqrt{E+1}$. Moreover
\[
 \frac{(g_2-E)!}{g_2!}\leq\left(\frac e{g_2}\right)^E.
\]
For example, integrate $\log t$ from $g_2-E$ to $g_2$ to bound
the logarithm of the falling factorial below by
$E(\log g_2-1)$. For sufficiently large $g$, $Ce/g_2\leq1/2$.
Thus the sum of \eqref{eq:s06:edge-mass} over $E\geq K$ is at most
\[
 C\mathcal T_g\left(\frac{Ce}{g_2}\right)^K
       \sum_{j\geq0}2^{-j}\sqrt{K+j+1}
 =O_K(\mathcal T_g g_2^{-K}).
\]
This bounds the complete tree sum, including all small-genus and
large-valence configurations.
\end{prf}

A one-edge tree has positive genus at each endpoint. Its two
orientations give the exact expression
\begin{equation}\label{eq:s06:one-edge}
 \frac12\sum_{G=1}^{g-1}\int_0^1
       W^{\rm top}_{G,1}(s)W^{\rm top}_{g-G,1}(s)\dd s.
\end{equation}
The factor $1/2$ also gives the correct automorphism weight when the
two genera agree. The next bound isolates the elliptic endpoints
without a fixed-split approximation.

\begin{lmm}[Nonelliptic one-edge tail]\label{lmm:s06:nonelliptic-tail}
The absolute contribution to \eqref{eq:s06:one-edge} of all splits
with $2\leq G\leq g-2$ is $O(\mathcal T_g g_2^{-3})$.
\end{lmm}
\begin{prf}
At an endpoint of genus $G\geq2$, the coefficient order is
$2G-2\geq2$. The two orders sum to $M=g_2-1$.
Applying the all-type vertex and beta bounds as above, using
$(m+1)(M-m+1)\geq M+1$, bounds the relative contribution by
\[
 C\sqrt{\frac{g_2+1}{M+1}}\,
       \frac{M!}{g_2!}
       \sum_{m=2}^{M-2}\binom{M}{m}^{-1}.
\]
Allowing odd $m$ only enlarges the sum. For $M\geq6$ it is at most
\[
 \frac2{\binom M2}+\frac{M-5}{\binom M3}=O(M^{-2}).
\]
Since $M!/g_2!=1/g_2$, the asserted bound follows. The finitely
many smaller totals are absorbed in the constant.
\end{prf}

\begin{exm}[Genus two]
Here $g_2=1$. The tree with no edges has value $\beta_2$, and the only
one-edge tree consists of two elliptic vertices with weight $1/2$.
The series \eqref{eq:s06:beta-series} gives
$\beta_1=1/24$ and $\beta_2=7/5760$. Thus
\[
 b_2=\frac{\beta_2+\beta_1^2/2}{12}=\frac1{5760}.
\]
This checks both the stack symmetry factor and the conversion by
$12^{g_2}$.
\end{exm}

% !TeX root = ../../main.tex
\subsection{The first correction and proof of Theorem 1.3}
\label{subsec:s06:04}

The tree bound leaves two calculations. The tree with no edges requires a
two-term coefficient expansion; an elliptic leaf needs only its
additive leading-amplitude version. We prove both versions on the
same Cauchy circle. In particular the error estimate does not require
division by an amplitude that might vanish at the saddle.

\begin{lmm}[Circular saddle]\label{lmm:s06:saddle}
Let $P$ be analytic on a neighborhood of the closed disk
$|x|\leq x_*$. As the positive integer $m\to\infty$,
\begin{equation}\label{eq:s06:additive-saddle}
 [x^m]\Phi(x)^{-m-1}P(x)
 =\frac{R^{-m}}{J_1(z_{0,1})\sqrt{2\pi m}}
             \left(P(x_*)+O_P(m^{-1})\right).
\end{equation}
The error is uniform for amplitudes analytic and uniformly bounded
on one fixed neighborhood of that disk. For $P=1$,
\begin{equation}\label{eq:s06:two-term-saddle}
 [x^m]\Phi(x)^{-m-1}
 =\frac{R^{-m}}{J_1(z_{0,1})\sqrt{2\pi m}}
 \left[1+\frac{-1/4+1/(24x_*)}{m}+O(m^{-2})\right].
\end{equation}
\end{lmm}
\begin{prf}
\emph{The integral and its outer arc.}
For $-\pi\leq\theta\leq\pi$ put $x_\theta=x_*e^{i\theta}$
and
$\mathcal A_P(\theta)=e^{i\theta}R P(x_\theta)/w(x_\theta)$.
Cauchy's formula is exactly
\begin{equation}\label{eq:s06:cauchy}
 [x^m]\Phi(x)^{-m-1}P(x)
 =\frac{R^{-m}}{2\pi\Phi(x_*)}
 \int_{-\pi}^{\pi}
 \left(\frac{w(x_\theta)}R\right)^{-m}
                         \mathcal A_P(\theta)\dd\theta.
\end{equation}
The analytic-disk hypothesis justifies this coefficient integral and
bounds its amplitude all around the circle. By
\eqref{eq:s06:circle-bound}, the part with
$|\theta|\geq m^{-2/5}$ is
$O_P(R^{-m}e^{-a m^{1/5}})$, after decreasing $a>0$ if needed.
It is negligible at every algebraic order relative to the leading
$R^{-m}/\sqrt m$ scale.

\emph{Local Taylor terms.}
Near zero define
$\mathcal H(\theta)=\log(w(x_\theta)/R)$ by the branch with
$\mathcal H(0)=0$. The series of $w$ satisfies $xw''+w=0$.
Together with $w(x_*)=R$ and $w'(x_*)=0$, differentiation gives
\begin{align*}
 \mathcal H(\theta)
 &=\frac{x_*}{2}\theta^2+\frac{i x_*}{3}\theta^3\\
 &\quad-\left(\frac{x_*^2}{12}+\frac{x_*}{8}\right)\theta^4
   -\frac{i x_*(7x_*+2)}{60}\theta^5+O(\theta^6),\\
 \mathcal A_1(\theta)
 &=1+i\theta-\frac{x_*+1}{2}\theta^2
                    -\frac{i(5x_*+1)}6\theta^3+O(\theta^4).
\end{align*}
Set $t=\sqrt m\,\theta$.
Define the three polynomials
\begin{align*}
 \mathcal Q_1(t)&=it-\frac{i x_*}{3}t^3,\\
 \mathcal Q_2(t)&=-\frac{x_*+1}{2}t^2
       +\left(\frac{x_*^2}{12}+\frac{11x_*}{24}\right)t^4
       -\frac{x_*^2}{18}t^6,\\
 \mathcal Q_3(t)&=i\biggl[-\frac{5x_*+1}{6}t^3
       +\left(\frac{11x_*^2}{30}+\frac{13x_*}{40}\right)t^5\\
 &\hspace{30mm}
       -\left(\frac{x_*^3}{36}+\frac{7x_*^2}{72}\right)t^7
       +\frac{x_*^3}{162}t^9\biggr].
\end{align*}

For $|t|\leq m^{1/10}$ the product of the Taylor expansions is
\begin{align}
 &e^{-m\mathcal H(t/\sqrt m)}\mathcal A_1(t/\sqrt m)\notag\\
 &\quad=e^{-x_*t^2/2}
 \left[1+\frac{\mathcal Q_1(t)}{\sqrt m}
         +\frac{\mathcal Q_2(t)}m
         +\frac{\mathcal Q_3(t)}{m^{3/2}}\right]\notag\\
 &\qquad\quad
   +O\left(m^{-2}(1+|t|^{12})e^{-x_*t^2/4}\right).
 \label{eq:s06:local-saddle}
\end{align}
For completeness, the Taylor remainders in the phase and amplitude
are $O(m^{-2}|t|^6)$ and $O(m^{-2}|t|^4)$ after rescaling.
The retained cubic exponential has size
$O(m^{-1/2}|t|^3)$. On this interval its ratio to $t^2$ tends
uniformly to zero, as do the higher phase terms. Exponential
Taylor's remainder is consequently bounded by $m^{-2}$ times a
polynomial of degree at most twelve, with the extra exponential
absorbed into $e^{-x_*t^2/4}$. This proves the integrable uniform
error in \eqref{eq:s06:local-saddle}.

\emph{Symmetry and the first correction.}
Both $\mathcal Q_1$ and $\mathcal Q_3$ are odd. Their integrals
against the even Gaussian over the symmetric interval vanish
exactly. The normalized Gaussian moments of $t^2,t^4,t^6$ are
$1/x_*$, $3/x_*^2$, $15/x_*^3$. They give
\[
 -\frac{x_*+1}{2x_*}
 +3\left(\frac1{12}+\frac{11}{24x_*}\right)
 -\frac{15}{18x_*}
 =-\frac14+\frac1{24x_*}.
\]
The Gaussian tails beyond $|t|=m^{1/10}$ are exponentially small,
and the error in \eqref{eq:s06:local-saddle} integrates to
$O(m^{-2})$. Finally
$\Phi(x_*)\sqrt{x_*}=J_1(z_{0,1})$ converts the prefactor in
\eqref{eq:s06:cauchy} to that in
\eqref{eq:s06:two-term-saddle}.

For general $P$, expand $\mathcal A_P$ through its quadratic term.
The term of order $m^{-1/2}$ is still odd; the remaining integral
is $P(x_*)$ plus an additive $O_P(m^{-1})$.
The same outer-arc bound applies. Cauchy bounds for derivatives on
a fixed neighborhood of the closed disk make these constants
uniform for the asserted families. This proves
\eqref{eq:s06:additive-saddle}, including when $P(x_*)=0$.
\end{prf}

The elliptic endpoint has an exact one-point value, so only its
opposite vertex must be estimated. The two placements of that
endpoint will remove the orientation half in the ordered split.

\begin{lmm}[Elliptic leaf]\label{lmm:s06:elliptic-leaf}
For $g\geq3$, the contribution to \eqref{eq:s06:one-edge} with
an elliptic endpoint is
\begin{equation}\label{eq:s06:elliptic-exact}
 \beta_1\beta_{g-1}(g_2-1)!
 [x^{g_2-1}]\Phi(x)^{-g_2}\Phi(x).
\end{equation}
Relative to
$\beta_g g_2!R^{-g_2}/(J_1(z_{0,1})\sqrt{2\pi g_2})$,
it equals
\begin{equation}\label{eq:s06:elliptic-relative}
 \frac{\pi^2J_1(z_{0,1})^2}{6g_2}+O(g_2^{-2}).
\end{equation}
\end{lmm}
\begin{prf}
The two ordered placements in \eqref{eq:s06:one-edge} cancel
its factor $1/2$. The small vertex is
$W^{\rm top}_{1,1}(s)=\beta_1=1/24$ by
\eqref{eq:s06:exact-vertex}. The entire series gives
\[
 \int_0^1J_0(2\sqrt{s x})\dd s
    =\sum_{a\geq0}\frac{(-x)^a}{a!(a+1)!}=\Phi(x).
\]
The other vertex has coefficient order $g_2-1$, proving
\eqref{eq:s06:elliptic-exact}.

Apply \eqref{eq:s06:additive-saddle} with $m=g_2-1$ and
$P=\Phi$. After division by the displayed scale its leading
coefficient is
\[
 \frac{\beta_1}{g_2}\frac{\beta_{g-1}}{\beta_g}
     R\Phi(x_*)\sqrt{\frac{g_2}{g_2-1}}.
\]
The omitted term is $O(g_2^{-2})$ in this ratio. Use
\eqref{eq:s06:beta-asymptotics} and
$R\Phi(x_*)=J_1(z_{0,1})^2$ to obtain
\eqref{eq:s06:elliptic-relative}.
\end{prf}

We can now combine a complete tree remainder with the two retained
contributions. The scale \eqref{eq:s06:tree-scale} differs by a
bounded positive factor from the leading saddle scale, so the
absolute tree estimates have exactly the relative strength needed
for the theorem.

\begin{prfn}{Theorem~\ref{thm:top-hodge}}
In \eqref{eq:s06:tree-expansion}, all trees with at least two edges
give $O(\mathcal T_g g_2^{-2})$ by
Proposition~\ref{prp:s06:tree-tail}. The nonelliptic one-edge sum
is $O(\mathcal T_g g_2^{-3})$ by
Lemma~\ref{lmm:s06:nonelliptic-tail}. The tree with no edges, by
\eqref{eq:s06:exact-vertex} and
\eqref{eq:s06:two-term-saddle} at coefficient order $g_2$, is
\[
 \frac{\beta_g g_2!R^{-g_2}}
      {J_1(z_{0,1})\sqrt{2\pi g_2}}
 \left[1+\frac{-1/4+1/(6z_{0,1}^2)}{g_2}
                 +O(g_2^{-2})\right].
\]
Lemma~\ref{lmm:s06:elliptic-leaf} supplies the remaining one-edge
term. Therefore
\begin{equation}\label{eq:s06:combined-expansion}
 \frac{12^{g_2}b_g}{g_2!}
 =\frac{\beta_g g_2!R^{-g_2}}
      {J_1(z_{0,1})\sqrt{2\pi g_2}}
       \left(1+\frac b{g_2}+O(g_2^{-2})\right).
\end{equation}
The $-1/4$ and $1/(6z_{0,1}^2)$ in the Introduction's scalar $b$
come from the saddle of the tree with no edges; its
$\pi^2J_1(z_{0,1})^2/6$ comes from the elliptic leaf.
Finally multiply by $g_2!/12^{g_2}$, use $12R=\rho$ and
\eqref{eq:s06:beta-asymptotics}. The exponentially small beta
error is absorbed into the remainder, giving
\[
 b_g=\frac1{J_1(z_{0,1})}\frac2{(2\pi)^{2g}}
 \frac{(g_2!)^2}{\rho^{g_2}\sqrt{2\pi g_2}}
       \left(1+\frac b{g_2}+O(g_2^{-2})\right).
\]
This is \eqref{eq:g111}.
\end{prfn}

